%! Rates of Change in Polar Functions — Unit 3, Lesson 3.15 — Homework
\documentclass[10pt]{article}
\usepackage{precalculus-article}
\usepackage{precalculus-boxes}
\newcommand{\TallMath}[1]{$\displaystyle #1\rule[-1.4em]{0pt}{3.2em}$}
\newcommand{\CourseName}{Precalculus}
\newcommand{\SchoolYear}{2026--2027}
\newcommand{\MeetingLength}{55 minutes}
\newcommand{\UnitNumberName}{Unit 3: Trigonometric and Polar Functions \quad}
\newcommand{\LessonNumberName}{Lesson 3.15: Rates of Change in Polar Functions}

\begin{document}

\pageheader{Unit 3, Lesson 3.15}{Homework}
\namedateperiod

\vspace{0.10in}

\begin{notesbox}{Problems 1--6}

\textbf{1.}\quad The table below gives values of $r = 3\sin\theta$ at key angles.
Classify each adjacent interval as P\&I (positive and increasing),
P\&D (positive and decreasing), N\&I (negative and increasing), or N\&D
(negative and decreasing).

\vspace{6pt}
\renewcommand{\arraystretch}{1.5}
\begin{tabular}{|c|c|c|c|c|c|}
\hline
$\theta$ & $0$ & $\pi/4$ & $\pi/2$ & $3\pi/4$ & $\pi$ \\
\hline
$r = 3\sin\theta$ & $0$ & $\approx 2.12$ & $3$ & $\approx 2.12$ & $0$ \\
\hline
\end{tabular}

\vspace{4pt}
\begin{tabular}{|c|c|c|}
\hline
\textbf{Interval} & \textbf{Classification} & \textbf{Curve behavior} \\
\hline
$[0, \pi/4]$ & \underline{\hspace{2.5cm}} & \underline{\hspace{3.5cm}} \\
\hline
$[\pi/4, \pi/2]$ & \underline{\hspace{2.5cm}} & \underline{\hspace{3.5cm}} \\
\hline
$[\pi/2, 3\pi/4]$ & \underline{\hspace{2.5cm}} & \underline{\hspace{3.5cm}} \\
\hline
$[3\pi/4, \pi]$ & \underline{\hspace{2.5cm}} & \underline{\hspace{3.5cm}} \\
\hline
\end{tabular}

\vspace{0.14in}
\textbf{2.}\quad Using the table from Problem 1, compute each average rate of change.
Show the formula and your work.

\vspace{6pt}
(a)\quad ARC of $r = 3\sin\theta$ from $\theta = 0$ to $\theta = \pi/4$:

\vspace{4pt}
ARC $= \dfrac{r(\pi/4) - r(0)}{\pi/4 - 0} = \dfrac{\rule{1.5cm}{0.4pt} - \rule{1.5cm}{0.4pt}}{\pi/4} = $ \blank{3cm}

\vspace{8pt}
(b)\quad ARC from $\theta = \pi/4$ to $\theta = \pi/2$:

\vspace{4pt}
ARC $= $ \blank{6cm}

\vspace{8pt}
(c)\quad ARC from $\theta = \pi/2$ to $\theta = \pi$:

\vspace{4pt}
ARC $= $ \blank{6cm}

\vspace{0.14in}
\textbf{3.}\quad The average rate of change of a polar function $r = g(\theta)$ over
an interval is $-1.5$.

\vspace{4pt}
(a)\quad What does this value tell you about whether $r$ is increasing or decreasing?

\writeline

\vspace{4pt}
(b)\quad If $r > 0$ throughout this interval, is the curve moving toward or away from the origin?

\writeline

\vspace{0.14in}
\textbf{4.}\quad A polar function $r = h(\theta)$ has $r = -0.8$ at $\theta = 2$ radians and
$r = -2.3$ at $\theta = 3$ radians.

\vspace{4pt}
(a)\quad Is $r$ negative and increasing (N\&I) or negative and decreasing (N\&D) on $[2, 3]$?

\vspace{4pt}
\underline{\hspace{5cm}}

\vspace{4pt}
(b)\quad In the context of a polar curve, is the traced point moving toward or away from the origin?
Explain.

\writeline

\vspace{4pt}
(c)\quad Compute the average rate of change of $r$ from $\theta = 2$ to $\theta = 3$.

\vspace{4pt}
ARC $= $ \blank{5cm}

\end{notesbox}

\vspace{0.08in}

\begin{notesbox}{Problems 5--6 (AP-Style Multiple Choice)}

\textbf{5.}\quad The table below shows selected values of a polar function $r = f(\theta)$.

\vspace{4pt}
\renewcommand{\arraystretch}{1.4}
\begin{tabular}{|c|c|c|c|c|}
\hline
$\theta$ & $0$ & $\pi/3$ & $2\pi/3$ & $\pi$ \\
\hline
$r = f(\theta)$ & $3$ & $1$ & $-1$ & $-3$ \\
\hline
\end{tabular}

\vspace{6pt}
Which of the following best describes the behavior of $r = f(\theta)$ on $[2\pi/3, \pi]$?

\begin{enumerate}[label=(\Alph*), leftmargin=2em, itemsep=3pt]
  \item Positive and increasing; the curve moves away from the origin.
  \item Positive and decreasing; the curve moves toward the origin.
  \item Negative and decreasing; the curve moves away from the origin (opposite ray).
  \item Negative and increasing; the curve moves toward the origin (opposite ray).
  \item Negative and decreasing; the curve moves toward the origin (opposite ray).
\end{enumerate}

Answer: \underline{\hspace{1.5cm}}

\vspace{0.12in}
\textbf{6.}\quad The average rate of change of $r = 2\sin\theta$ over $[\pi/4, \pi/2]$ is
closest to which of the following?

\begin{enumerate}[label=(\Alph*), leftmargin=2em, itemsep=3pt]
  \item $-1.08$
  \item $0$
  \item $0.54$
  \item $1.08$
  \item $2.55$
\end{enumerate}

Answer: \underline{\hspace{1.5cm}}

\end{notesbox}

\vspace{0.08in}

\begin{tcolorbox}[colback=navylight, colframe=navy, arc=2mm,
    left=3mm, right=3mm, top=2mm, bottom=2mm]
\small\textbf{Preview --- Unit 4: Functions Involving Parameters, Vectors, and Matrices}\\
This is the final homework of Unit 3. In Unit 4 you will study parametric functions,
where $x$ and $y$ both depend on a parameter $t$. The average rate of change of $x$
(or $y$) with respect to $t$ will tell you direction and speed along the parametric curve ---
a direct extension of the rate-of-change ideas you used today for polar functions.
\end{tcolorbox}

\end{document}
